\section{Discussion}
\label{sec:discussion}

\para{Dilution, not saturation.}
H4 pictures rule-following as a capacity that fills up: a model handles rules perfectly until it runs out of room.
Our data do not show this.
Every model loses compliance from about the fifth rule onward, and the curve is smooth rather than a cliff.
A better description is dilution: each additional rule, and each additional thousand tokens between a rule and the user turn, lowers the chance that any given rule shapes the output.
Model quality sets both the starting level and the rate of decline, but no tested model stays flat to 20 rules.
\qwenthree comes closest (0.92 at 10 rules, 0.87 at 20), which suggests that newer instruction-tuning recipes matter more than raw scale.
\qfourteen did not beat \qseven; this could be partly due to quantization, but it fits the view that scale alone does not buy rule capacity.

\para{Length hurts through distance and burial.}
Comparing \expB and \expF, the same five rules in a 4k-token prompt are followed 28\% of the time when scattered, 34\% as a block in the middle, 40\% as a block at the top, and 55\% as a block at the bottom.
Chat models appear to weight the most recent system text heavily when they decide on format and content.
For short lists, the first rules benefit instead.
The two patterns are consistent once we note that in a short list all rules are already close to the end.

\para{Contradictions are costly.}
Contradictory rules do more than force a choice.
For numeric constraints, models often split the difference and fail both rules, so a contradictory pair can be worse than dropping one of its rules.
For soft conflicts, the ``thorough'' instruction reliably wins.

\para{Practical guidance.}
Our results suggest three rules of thumb for prompt authors.
First, every rule has a cost, starting with the first few, so prompts should contain only rules that matter.
Second, in long system prompts, critical rules should sit in one contiguous block at the end, right before the user turn; in short rule lists, the most important rules should come first.
Third, prompts should be checked for contradictions, because models seldom resolve them gracefully.

\subsection{Limitations}
\label{sec:limitations}

\para{Open models only.}
No API keys were available, so we could not test commercial frontier models.
These may be much more robust; \qwenthree already is.

\para{Quantization is confounded with size.}
The 3B, 8B, and 14B models are AWQ int4, while the 0.5B--1.5B and 7B models are bf16.
The finding that \qfourteen does not beat \qseven may partly reflect quantization.

\para{Incomplete coverage.}
Shared-GPU contention meant that \expB covers 7 models, \expE and \expF cover 6, and \qwenthree and \qfourteen have only \expA and \expD.
Our length and placement conclusions therefore do not include the two strongest models on the count experiment.

\para{Synthetic rules and padding.}
Our rules are verifiable but artificial: arbitrary keywords and formats.
Real system-prompt rules (tone, refusals, tool policies) may behave differently.
Inclusion rules are harder than natural ones, and exclusion rules sit near ceiling.
Our padding is encyclopedic text labeled as background; realistic padding such as tool documentation may compete with rules more or less strongly.

\para{Design confounds and variance.}
In \expA, count and length grow together; \expE addresses this only for the models that ran it.
Each model was run once with greedy decoding, and run-to-run variance from batch non-invariance was checked on one model only.
Our confidence intervals cover query sampling, not decoding noise.
Finally, we test only single-turn answers of about 100--500 words.

\para{Broader impact.}
System-prompt rules often encode safety and policy constraints.
Our results show that such rules can be silently dropped in long prompts, so deployers should not assume that a rule written into a system prompt will be enforced.
Measuring rule compliance under realistic prompt lengths should be part of deployment testing.
